% ============================================================
\chapter{Le Processus de Mod\'elisation}
\label{ch:processus}
% ============================================================

Galil\'ee \'ecrivait que le livre de la nature est \'ecrit en langage
math\'ematique. Mais entre la nature et les \'equations, il y a un pont \`a
construire~: c'est le \emph{mod\`ele}. Un mod\`ele est une simplification
d\'elib\'er\'ee de la r\'ealit\'e, qui capture les m\'ecanismes essentiels tout en
ignorant les d\'etails superflus. L'art de la mod\'elisation --- car c'est
bien un art autant qu'une science --- consiste \`a choisir le bon niveau
d'abstraction~: assez de d\'etails pour que le mod\`ele soit utile, assez de
simplicit\'e pour qu'il soit traitable.

\begin{center}
\textit{<<\,La science est faite de donn\'ees comme une maison est faite
de pierres ; mais une accumulation de donn\'ees n'est pas plus une
science qu'un tas de pierres n'est une maison.\,>>}\\[4pt]
--- Henri Poincar\'e
\end{center}

\bigskip

Ce premier chapitre pr\'esente la d\'emarche g\'en\'erale de mod\'elisation
math\'ematique.  Nous d\'ecrivons les \'etapes cl\'es, de l'observation du
ph\'enom\`ene \`a la validation du mod\`ele, et nous classifions les grandes
familles de mod\`eles math\'ematiques.

% ------------------------------------------------------------
\section{Qu'est-ce qu'un mod\`ele math\'ematique\,?}
% ------------------------------------------------------------

\begin{definition}[Mod\`ele math\'ematique]
Un \textbf{mod\`ele math\'ematique} est une repr\'esentation formelle d'un
syst\`eme r\'eel (physique, biologique, \'economique, etc.) \`a l'aide
d'objets math\'ematiques (\'equations, in\'egalit\'es, graphes, distributions
de probabilit\'e) qui capturent les relations essentielles entre les
grandeurs du syst\`eme.
\end{definition}

\begin{remarque}
Un mod\`ele n'est jamais unique : le m\^eme ph\'enom\`ene peut \^etre
d\'ecrit par diff\'erents mod\`eles, de niveaux de complexit\'e vari\'es.
Le choix du mod\`ele d\'epend de la question pos\'ee et des donn\'ees
disponibles.
\end{remarque}

% ------------------------------------------------------------
\section{Les \'etapes du processus de mod\'elisation}
\label{sec:etapes}
% ------------------------------------------------------------

Le processus de mod\'elisation se d\'ecompose en cinq grandes \'etapes,
formant un cycle it\'eratif.

\begin{enumerate}
  \item \textbf{Observation et identification du probl\`eme.}
    On d\'elimite le ph\'enom\`ene \`a \'etudier, on identifie les grandeurs
    pertinentes (variables d'\'etat, param\`etres, entr\'ees/sorties)
    et on rassemble les donn\'ees disponibles.

  \item \textbf{Formulation des hypoth\`eses.}
    On \'enonce explicitement les simplifications adopt\'ees :
    quelles interactions sont n\'eglig\'ees, quelles grandeurs sont
    suppos\'ees constantes, quel est le domaine de validit\'e spatial
    et temporel.

  \item \textbf{Construction du mod\`ele math\'ematique.}
    On traduit les hypoth\`eses en \'equations.  Cela peut donner :
    \begin{itemize}
      \item une \'equation diff\'erentielle ordinaire (EDO),
      \item une \'equation aux d\'eriv\'ees partielles (EDP),
      \item un syst\`eme d'\'equations alg\'ebriques,
      \item un programme d'optimisation,
      \item un mod\`ele stochastique.
    \end{itemize}

  \item \textbf{R\'esolution et analyse.}
    On r\'esout le mod\`ele analytiquement (quand c'est possible) ou
    num\'eriquement.  On \'etudie le comportement qualitatif :
    existence et unicit\'e des solutions, stabilit\'e des \'equilibres,
    d\'ependance aux param\`etres.

  \item \textbf{Validation et interpr\'etation.}
    On compare les pr\'edictions du mod\`ele aux donn\'ees exp\'erimentales.
    Si l'accord est insuffisant, on revient \`a l'\'etape~2 pour modifier
    les hypoth\`eses.
\end{enumerate}

\begin{intuition}[title=\textbf{Le cycle de mod\'elisation}]
La mod\'elisation n'est pas un processus lin\'eaire mais un \textbf{cycle}.
Apr\`es la validation, on revient souvent raffiner les hypoth\`eses,
enrichir le mod\`ele ou en r\'eduire la complexit\'e.  Un bon mod\'elisateur
effectue plusieurs tours de ce cycle.
\end{intuition}

% ------------------------------------------------------------
\section{Classification des mod\`eles}
\label{sec:classification}
% ------------------------------------------------------------

\begin{definition}[Types de mod\`eles]
On distingue les mod\`eles selon plusieurs crit\`eres :
\begin{enumerate}
  \item \textbf{D\'eterministe vs.\ stochastique.}
    Un mod\`ele d\'eterministe pr\'edit une trajectoire unique pour des
    conditions initiales donn\'ees ; un mod\`ele stochastique incorpore
    de l'al\'eatoire.
  \item \textbf{Continu vs.\ discret.}
    Le temps et/ou l'espace peuvent varier de fa\c{c}on continue ou discr\`ete.
  \item \textbf{Statique vs.\ dynamique.}
    Un mod\`ele statique d\'ecrit un \'etat d'\'equilibre ; un mod\`ele
    dynamique d\'ecrit l'\'evolution temporelle.
  \item \textbf{Lin\'eaire vs.\ non lin\'eaire.}
    La lin\'earit\'e des relations entre variables conditionne fortement
    les m\'ethodes d'analyse.
  \item \textbf{Param\'etrique vs.\ non param\'etrique.}
    Selon que la structure du mod\`ele est fix\'ee a priori ou
    d\'eduite des donn\'ees.
\end{enumerate}
\end{definition}

\begin{exemple}
Consid\'erons la chute d'un objet. En n\'egligeant la r\'esistance de
l'air, le mod\`ele est :
\[
  m\ddot{y} = -mg, \qquad y(0)=h,\quad \dot{y}(0)=0.
\]
C'est un mod\`ele d\'eterministe, continu, dynamique et lin\'eaire.
Sa solution est $y(t)=h-\tfrac{1}{2}gt^2$.

Si l'on ajoute la r\'esistance de l'air proportionnelle \`a $\dot{y}^2$,
on obtient un mod\`ele \textbf{non lin\'eaire} :
\[
  m\ddot{y} = -mg + k\dot{y}^2.
\]
\end{exemple}

% ------------------------------------------------------------
\section{Principes fondamentaux de construction}
\label{sec:principes}
% ------------------------------------------------------------

\subsection{Lois de conservation}

De nombreux mod\`eles reposent sur des \textbf{lois de conservation} :
conservation de la masse, de l'\'energie, de la quantit\'e de mouvement,
ou du nombre d'individus.

\begin{theoreme}[Bilan g\'en\'erique]
Pour une quantit\'e $Q(t)$ contenue dans un domaine $\Omega$, la loi
de conservation g\'en\'erique s'\'ecrit :
\[
  \frac{dQ}{dt} = \text{(flux entrant)} - \text{(flux sortant)}
                 + \text{(source)} - \text{(puits)}.
\]
\end{theoreme}

\begin{exemple}
Soit $N(t)$ la taille d'une population. Si le taux de naissance
est $b$ et le taux de mortalit\'e est $d$ (par individu par unit\'e
de temps), la loi de conservation donne :
\[
  \frac{dN}{dt} = bN - dN = (b-d)N.
\]
C'est le mod\`ele de Malthus, que nous \'etudierons en d\'etail au
chapitre~\ref{ch:populations}.
\end{exemple}

\subsection{Lois constitutives}

Les lois constitutives relient les flux aux grandeurs d'\'etat.

\begin{exemple}
\begin{itemize}
  \item \textbf{Loi de Fourier :} le flux de chaleur est proportionnel
    au gradient de temp\'erature : $\mathbf{q} = -\kappa\nabla T$.
  \item \textbf{Loi de Fick :} le flux de diffusion est proportionnel
    au gradient de concentration : $\mathbf{J} = -D\nabla c$.
  \item \textbf{Loi de Hooke :} la contrainte est proportionnelle
    \`a la d\'eformation : $\sigma = E\varepsilon$.
\end{itemize}
\end{exemple}

\subsection{Principes variationnels}

Certains mod\`eles sont obtenus en minimisant (ou en rendant stationnaire)
une fonctionnelle.

\begin{definition}[Fonctionnelle d'action]
En m\'ecanique, le principe de moindre action \'enonce que la trajectoire
r\'eelle rend stationnaire la fonctionnelle
\[
  \mathcal{S}[q] = \int_{t_1}^{t_2} L(q,\dot{q},t)\,dt,
\]
o\`u $L$ est le lagrangien du syst\`eme.  Les \'equations d'Euler--Lagrange
qui en d\'ecoulent seront \'etudi\'ees au chapitre~\ref{ch:optimisation}.
\end{definition}

% ------------------------------------------------------------
\section{Qualit\'es d'un bon mod\`ele}
% ------------------------------------------------------------

\begin{bonnepratique}[title=\textbf{Crit\`eres de qualit\'e}]
Un bon mod\`ele satisfait un compromis entre :
\begin{itemize}
  \item \textbf{Simplicit\'e} (parcimonie) : le nombre de param\`etres
    doit rester raisonnable (principe du rasoir d'Occam).
  \item \textbf{Pr\'ecision} : les pr\'edictions doivent \^etre
    suffisamment proches des observations.
  \item \textbf{Robustesse} : de petites variations des param\`etres
    ne doivent pas entra\^iner de changements dramatiques des pr\'edictions.
  \item \textbf{Interpr\'etabilit\'e} : les param\`etres doivent avoir
    une signification physique ou biologique.
\end{itemize}
\end{bonnepratique}

\begin{proposition}[Principe de parcimonie]
\'Etant donn\'es deux mod\`eles de pr\'ecision comparable, on pr\'ef\`ere
celui qui poss\`ede le moins de param\`etres libres.  Ce principe est
formalis\'e par des crit\`eres d'information tels que l'AIC (Akaike
Information Criterion) :
\[
  \mathrm{AIC} = 2k - 2\ln \hat{L},
\]
o\`u $k$ est le nombre de param\`etres et $\hat{L}$ la vraisemblance
maximale.
\end{proposition}

% ------------------------------------------------------------
\section{Sensibilit\'e et incertitude}
\label{sec:sensibilite}
% ------------------------------------------------------------

\begin{definition}[Analyse de sensibilit\'e]
L'analyse de sensibilit\'e \'etudie comment la sortie $y$ d'un mod\`ele
varie lorsqu'un param\`etre $p$ change.  L'\textbf{indice de sensibilit\'e
locale} est d\'efini par :
\[
  S_p = \frac{p}{y}\,\frac{\partial y}{\partial p}.
\]
Un indice $\abs{S_p}\gg 1$ signale que le mod\`ele est tr\`es sensible
au param\`etre $p$.
\end{definition}

\begin{exemple}
Pour le mod\`ele exponentiel $y(t)=y_0 e^{rt}$, on a
\[
  S_r = \frac{r}{y}\,\frac{\partial y}{\partial r}
      = \frac{r}{y_0 e^{rt}}\,y_0 t e^{rt} = rt.
\]
La sensibilit\'e au taux $r$ cro\^it lin\'eairement avec le temps :
de petites erreurs sur $r$ se traduisent par de grandes erreurs
sur $y$ \`a long terme.
\end{exemple}

\begin{codeblock}[title=\textbf{Analyse de sensibilit\'e en Python}]
\begin{minted}{python}
import numpy as np
import matplotlib.pyplot as plt

def model(t, r, y0=1.0):
    return y0 * np.exp(r * t)

t = np.linspace(0, 10, 200)
r_nom = 0.5
dr = 0.05  # perturbation de 10%

y_nom = model(t, r_nom)
y_plus = model(t, r_nom + dr)
y_minus = model(t, r_nom - dr)

fig, (ax1, ax2) = plt.subplots(1, 2, figsize=(12, 4))

ax1.plot(t, y_nom, 'b-', label=f'$r = {r_nom}$')
ax1.fill_between(t, y_minus, y_plus, alpha=0.3,
                 label=f'$r = {r_nom} \\pm {dr}$')
ax1.set_xlabel('$t$'); ax1.set_ylabel('$y(t)$')
ax1.legend(); ax1.set_title('Solutions')

S_r = r_nom * t  # indice de sensibilité
ax2.plot(t, S_r, 'r-')
ax2.set_xlabel('$t$'); ax2.set_ylabel('$S_r$')
ax2.set_title('Indice de sensibilité $S_r = rt$')
plt.tight_layout(); plt.savefig('sensibilite.pdf')
\end{minted}
\end{codeblock}

% ------------------------------------------------------------
\section{Ajustement aux donn\'ees}
\label{sec:ajustement}
% ------------------------------------------------------------

Une fois le mod\`ele construit, il faut en estimer les param\`etres
\`a partir des donn\'ees.

\begin{definition}[Moindres carr\'es]
\'Etant donn\'ees des observations $(t_i, y_i^{\mathrm{obs}})$,
$i=1,\ldots,n$, et un mod\`ele $y(t;\mathbf{p})$ d\'ependant d'un
vecteur de param\`etres $\mathbf{p}$, l'estimateur des
\textbf{moindres carr\'es} est :
\[
  \hat{\mathbf{p}} = \arg\min_{\mathbf{p}}
    \sum_{i=1}^{n} \bigl(y_i^{\mathrm{obs}} - y(t_i;\mathbf{p})\bigr)^2.
\]
\end{definition}

\begin{codeblock}[title=\textbf{Ajustement par moindres carr\'es}]
\begin{minted}{python}
import numpy as np
from scipy.optimize import curve_fit

# Données synthétiques : croissance exponentielle bruitée
np.random.seed(42)
t_data = np.linspace(0, 5, 20)
r_true, y0_true = 0.8, 2.0
y_data = y0_true * np.exp(r_true * t_data) \
         + np.random.normal(0, 1, len(t_data))

def modele(t, y0, r):
    return y0 * np.exp(r * t)

popt, pcov = curve_fit(modele, t_data, y_data, p0=[1, 0.5])
print(f"Paramètres estimés : y0 = {popt[0]:.3f}, r = {popt[1]:.3f}")
print(f"Écarts-types : {np.sqrt(np.diag(pcov))}")
\end{minted}
\end{codeblock}

% ------------------------------------------------------------
\section{Exemple complet : refroidissement d'une tasse de caf\'e}
\label{sec:cafe}
% ------------------------------------------------------------

Illustrons le processus complet sur un exemple classique.

\paragraph{\'Etape 1 : Observation.}
Une tasse de caf\'e chaude pos\'ee sur une table refroidit progressivement.
On mesure la temp\'erature $T(t)$ au cours du temps.

\paragraph{\'Etape 2 : Hypoth\`eses.}
\begin{itemize}
  \item La temp\'erature du caf\'e est homog\`ene (mod\`ele \`a
    param\`etres group\'es).
  \item La temp\'erature ambiante $T_a$ est constante.
  \item Le flux de chaleur est proportionnel \`a l'\'ecart de
    temp\'erature (loi de Newton du refroidissement).
\end{itemize}

\paragraph{\'Etape 3 : Mod\`ele.}
La loi de Newton donne :
\[
  \frac{dT}{dt} = -k(T - T_a), \qquad T(0)=T_0,
\]
o\`u $k>0$ est le coefficient de transfert thermique.

\paragraph{\'Etape 4 : R\'esolution.}
L'EDO est lin\'eaire du premier ordre.  Sa solution est :
\[
  T(t) = T_a + (T_0 - T_a)\,e^{-kt}.
\]

\begin{formules}[title=\textbf{Loi de Newton du refroidissement}]
\[
  T(t) = T_a + (T_0 - T_a)\,e^{-kt}, \qquad k > 0.
\]
\begin{itemize}
  \item $T_0$ : temp\'erature initiale.
  \item $T_a$ : temp\'erature ambiante.
  \item $k$ : constante de refroidissement (en $\text{s}^{-1}$).
  \item Temps de demi-refroidissement : $t_{1/2} = \ln 2 / k$.
\end{itemize}
\end{formules}

\paragraph{\'Etape 5 : Validation.}

\begin{codeblock}[title=\textbf{Simulation et ajustement --- refroidissement}]
\begin{minted}{python}
import numpy as np
import matplotlib.pyplot as plt
from scipy.optimize import curve_fit

# Données expérimentales (simulées)
np.random.seed(0)
t_exp = np.arange(0, 31, 2)  # minutes
T_a, T0_true, k_true = 22.0, 90.0, 0.07
T_exp = T_a + (T0_true - T_a) * np.exp(-k_true * t_exp) \
        + np.random.normal(0, 1.5, len(t_exp))

def newton_cooling(t, T0, k):
    return T_a + (T0 - T_a) * np.exp(-k * t)

popt, pcov = curve_fit(newton_cooling, t_exp, T_exp, p0=[85, 0.05])
print(f"T0 = {popt[0]:.1f} °C,  k = {popt[1]:.4f} /min")

t_fine = np.linspace(0, 30, 300)
plt.figure(figsize=(8, 4))
plt.scatter(t_exp, T_exp, c='red', label='Données')
plt.plot(t_fine, newton_cooling(t_fine, *popt), 'b-',
         label='Modèle ajusté')
plt.axhline(T_a, color='gray', ls='--', label=f'$T_a = {T_a}$ °C')
plt.xlabel('Temps (min)'); plt.ylabel('Température (°C)')
plt.legend(); plt.title('Refroidissement de Newton')
plt.tight_layout(); plt.savefig('refroidissement.pdf')
\end{minted}
\end{codeblock}

% ------------------------------------------------------------
\section{Limites et pi\`eges courants}
% ------------------------------------------------------------

\begin{attention}[title=\textbf{Pi\`eges \`a \'eviter}]
\begin{enumerate}
  \item \textbf{Sur-ajustement (overfitting).}
    Un mod\`ele avec trop de param\`etres peut s'ajuster parfaitement
    aux donn\'ees d'entra\^inement mais \'echouer en pr\'ediction.
  \item \textbf{Extrapolation abusive.}
    Un mod\`ele valid\'e sur un intervalle $[0,T]$ ne doit pas \^etre
    extrapol\'e sans pr\'ecaution \`a $t\gg T$.
  \item \textbf{Confusion corr\'elation/causalit\'e.}
    Un bon ajustement ne prouve pas un lien causal.
  \item \textbf{Oubli des unit\'es.}
    Ne jamais additionner des grandeurs de dimensions diff\'erentes
    (voir chapitre~\ref{ch:dimensionnelle}).
\end{enumerate}
\end{attention}

% ------------------------------------------------------------
\section{Exercices}
% ------------------------------------------------------------

\begin{exercice}
Un r\'eservoir contient initialement $V_0=1000$~L d'eau pure. On y
verse une solution sal\'ee de concentration $c_{\mathrm{in}}=30$~g/L
au d\'ebit $q_{\mathrm{in}}=5$~L/min.  Le m\'elange, suppos\'e homog\`ene,
sort au d\'ebit $q_{\mathrm{out}}=5$~L/min.

\begin{enumerate}
  \item \'Ecrire l'EDO v\'erifi\'ee par la quantit\'e de sel $Q(t)$ dans
    le r\'eservoir.
  \item R\'esoudre l'EDO et d\'eterminer la concentration \`a l'\'equilibre.
  \item Tracer $Q(t)$ pour $t\in[0,400]$~min en Python.
  \item Discuter ce qui change si $q_{\mathrm{out}}\neq q_{\mathrm{in}}$.
\end{enumerate}
\end{exercice}

\begin{exercice}
On consid\`ere le mod\`ele $y(t)=A\sin(\omega t+\varphi)$ pour
d\'ecrire un signal p\'eriodique.
\begin{enumerate}
  \item G\'en\'erer des donn\'ees synth\'etiques bruit\'ees avec
    $A=3$, $\omega=2\pi/5$, $\varphi=\pi/4$ et un bruit gaussien
    d'\'ecart-type~$0{,}5$.
  \item Ajuster le mod\`ele par moindres carr\'es.
  \item Calculer les intervalles de confiance \`a 95\,\% pour
    $A$, $\omega$ et $\varphi$.
\end{enumerate}
\end{exercice}

\begin{exercice}
Pour le mod\`ele de refroidissement de Newton :
\begin{enumerate}
  \item Montrer que la temp\'erature $T(t)$ est strictement d\'ecroissante
    si $T_0>T_a$ et strictement croissante si $T_0<T_a$.
  \item Calculer l'indice de sensibilit\'e $S_k$ et interpr\'eter.
  \item Modifier le mod\`ele pour tenir compte du rayonnement
    (loi de Stefan--Boltzmann : flux $\propto T^4 - T_a^4$).
    R\'esoudre num\'eriquement et comparer.
\end{enumerate}
\end{exercice}

\begin{exercice}
On observe la d\'ecroissance d'une substance radioactive.  Le nombre
d'atomes $N(t)$ satisfait $dN/dt = -\lambda N$.
\begin{enumerate}
  \item Montrer que la demi-vie est $t_{1/2}=\ln 2/\lambda$.
  \item \'Ecrire un programme Python qui, \`a partir de donn\'ees
    $(t_i,N_i)$, estime $\lambda$ par moindres carr\'es et
    par r\'egression lin\'eaire sur $\ln N$.
  \item Comparer les deux m\'ethodes et discuter leurs avantages.
\end{enumerate}
\end{exercice}
